\begingroup
% Kiểu cột L cho tabularx (khai báo trong nhóm riêng, theo quy ước của báo cáo).
\newcolumntype{L}[1]{>{\hsize=#1\hsize\raggedright\arraybackslash}X}

\section{Hướng dẫn cấu hình bộ định tuyến và các thiết bị}
\label{sec:configuration-guide}

Phần này ghi lại quy trình cấu hình đã thực hiện trên thiết bị thật của nhóm, kèm
các giá trị thực tế đọc được trên giao diện quản trị. Tên menu tương ứng với bản
OpenWrt cài trên bộ định tuyến; với phần sụn của nhà sản xuất khác, tên mục có thể khác
nhưng nguyên tắc cấu hình không đổi. Không đăng tải mật khẩu quản trị hoặc mật
khẩu Wi-Fi thật vào báo cáo.

\subsection{Thiết bị và thông số thực tế}
\label{subsec:thiet-bi-thuc-te}

\begin{table}[htbp]
  \centering\footnotesize
  \caption{Thông số thiết bị và cấu hình mạng đã thiết lập trên N09-Router.}
  \label{tab:cauhinh-thucte}
  \begin{tabularx}{\linewidth}{@{} L{0.62} L{1.38} @{}}
    \toprule
    \textbf{Hạng mục} & \textbf{Giá trị thực tế} \\
    \midrule
    Mẫu thiết bị / phần mềm & TP-Link Archer C20 v5, phần sụn OpenWrt, giao diện quản trị
      LuCI \\
    Tên thiết bị & \texttt{N09-Router} \\
    Địa chỉ quản trị & \texttt{192.168.1.1} (giao diện \texttt{br-lan}),
      \texttt{192.168.1.0/24} \\
    MAC giao diện LAN & \texttt{B0:95:75:9F:1C:A3} \\
    Bộ phát 2,4~GHz & MediaTek MT7628, 802.11b/g/n, kênh~6 (2,437~GHz), rộng
      20~MHz, mã quốc gia VN \\
    \ac{ssid} 2,4~GHz & \texttt{25CLC04-N09}, WPA2-PSK, mã hóa CCMP (AES) \\
    Bộ phát 5~GHz & MediaTek MT7610E, 802.11ac, kênh~36 (5,180~GHz), rộng
      80~MHz \\
    \ac{ssid} 5~GHz & \texttt{25CLC04-N09-5G}, WPA2-PSK, mã hóa CCMP; \ac{bssid}
      \texttt{B0:95:75:9F:1C:A2} \\
    Mạng khách & \texttt{25CLC04-N09-Guest}, dải \texttt{192.168.20.0/24}, cổng
      \texttt{192.168.20.1}, DHCP riêng, vùng tường lửa riêng \\
    Cấp địa chỉ & DHCP trên \ac{lan}, thời hạn thuê 12~h; có đặt IP tĩnh theo MAC
      cho máy bàn (\texttt{192.168.1.50}) \\
    Cổng WAN & Chế độ DHCP client; trong sơ đồ demo cuối cùng cổng WAN không nối
      ra Internet \\
    \bottomrule
  \end{tabularx}
\end{table}

\subsection{Cấu hình cơ bản trên giao diện LuCI}
\begin{enumerate}[leftmargin=*]
  \item \textbf{Đăng nhập và đặt mật khẩu quản trị.} Máy quản trị nối vào cổng
    LAN (hoặc vào Wi-Fi tạm), mở \texttt{http://192.168.1.1}. Router mới ở trạng
    thái chưa có mật khẩu nên hiện cảnh báo; đặt mật khẩu cho tài khoản
    \texttt{root} ngay từ lần đăng nhập đầu tiên. Đây là bước bắt buộc vì ai cắm
    được dây vào cổng LAN cũng có thể thay đổi toàn bộ cấu hình.
  \item \textbf{Đặt tên thiết bị và múi giờ.} Vào \emph{System $\rightarrow$
    System}, đổi Hostname thành \texttt{N09-Router}, chọn múi giờ
    \emph{Asia/Ho\_Chi\_Minh} rồi đồng bộ giờ với máy quản trị. Việc đặt đúng
    múi giờ là điều kiện để các luật chặn theo khung giờ ở
    Mục~\ref{subsec:cauhinh-nangcao} hoạt động đúng.
  \item \textbf{Cấu hình \ac{lan} và DHCP.} Vào \emph{Network $\rightarrow$
    Interfaces}, đặt địa chỉ \ac{lan} \texttt{192.168.1.1/24} và bật máy chủ DHCP
    cho dải này. Cổng WAN để ở chế độ DHCP client để nhận địa chỉ từ modem hoặc
    router phía trên khi cần nối Internet.
  \item \textbf{Cấu hình Wi-Fi 2,4~GHz.} Vào \emph{Network $\rightarrow$
    Wireless}, mở thiết bị MT7628: chọn Country Code \texttt{VN}, Mode \texttt{N},
    Channel \texttt{6}, Width \texttt{20~MHz}; phần Interface đặt ESSID
    \texttt{25CLC04-N09}, Mode \emph{Access Point}, Network \texttt{lan}.
  \item \textbf{Cấu hình Wi-Fi 5~GHz.} Mở thiết bị MT7610E: Country Code
    \texttt{VN}, Mode \texttt{AC}, Channel \texttt{36}, Width \texttt{80~MHz};
    ESSID \texttt{25CLC04-N09-5G}.
  \item \textbf{Cấu hình bảo mật.} Ở tab \emph{Wireless Security} của mỗi mạng,
    chọn Encryption WPA2-PSK, Cipher \emph{Force CCMP (AES)} và đặt khóa dài, khó
    đoán. Sau khi bấm \emph{Save \& Apply}, trang Wireless hiển thị hai mạng đang
    phát, mỗi mạng kèm một \ac{bssid} riêng -- đúng với mô tả lý thuyết về
    \ac{bss} và \ac{bssid}.
\end{enumerate}

\subsection{Các chức năng quản lý đã cấu hình và thử}
\label{subsec:cauhinh-nangcao}

Các chức năng dưới đây được thiết lập trên N09-Router và được quay lại trong
video minh chứng (Mục~\ref{sec:video}). Cột ``Kết quả quan sát'' chỉ ghi những gì
thực sự thấy trên màn hình, không suy diễn thêm.

\begin{table}[htbp]
  \centering\footnotesize
  \caption{Các chức năng quản lý đã cấu hình trên router và kết quả quan sát.}
  \label{tab:chucnang-router}
  \begin{tabularx}{\linewidth}{@{} L{0.7} L{1.1} L{1.2} @{}}
    \toprule
    \textbf{Chức năng} & \textbf{Đường dẫn cấu hình} & \textbf{Kết quả quan sát} \\
    \midrule
    Ẩn tên mạng & Wireless $\rightarrow$ Hide ESSID & Mạng 2,4~GHz biến mất khỏi
      danh sách quét của thiết bị thử, mạng 5~GHz vẫn hiện; bật lại thì mạng hiện
      lại bình thường \\
    Mạng khách tách riêng & Interfaces + Wireless + Firewall zone
      \texttt{guest} & Thiết bị thử nhận địa chỉ trong dải
      \texttt{192.168.20.0/24}, ra Internet bình thường nhưng không mở được trang
      quản trị \texttt{192.168.1.1} \\
    Lọc địa chỉ MAC & Wireless $\rightarrow$ MAC-Filter & Thêm MAC thiết bị thử
      vào danh sách chặn: thiết bị bị ngắt và không kết nối lại được dù đúng mật
      khẩu; tắt lọc thì vào lại được \\
    Cấp IP tĩnh theo MAC & DHCP and DNS $\rightarrow$ Static Leases & Máy bàn tắt
      rồi bật Wi-Fi, nhận đúng \texttt{192.168.1.50} như đã gán \\
    Chặn tên miền & DHCP and DNS $\rightarrow$ DNS Records & Trỏ
      \texttt{vnexpress.net} và \texttt{www.vnexpress.net} về \texttt{0.0.0.0}:
      trình duyệt không mở được trang này, các trang khác vẫn bình thường; xóa
      luật thì truy cập lại được \\
    Giới hạn giờ truy cập & Firewall $\rightarrow$ Traffic Rules + Time
      Restrictions & Luật chặn theo MAC nguồn ra vùng \texttt{wan}: thiết bị vẫn
      báo có Wi-Fi nhưng không mở được trang web; tắt luật thì vào lại được \\
    Chuyển tiếp cổng & Firewall $\rightarrow$ Port Forwards & Từ ngoài mạng, mở
      địa chỉ WAN cổng 8080 truy cập được dịch vụ trên máy bàn; xóa luật thì
      không còn truy cập \\
    NAT và chẩn đoán & Firewall (Masquerading), Network $\rightarrow$
      Diagnostics & Vùng \texttt{wan} bật Masquerading (NAT); Traceroute tới
      \texttt{8.8.8.8} cho thấy gói tin đi qua router của mạng nhà trước khi ra
      nhà cung cấp dịch vụ \\
    Sao lưu cấu hình & System $\rightarrow$ Backup / Flash Firmware & Tạo được
      tệp lưu trữ cấu hình để phục hồi khi cần \\
    \bottomrule
  \end{tabularx}
\end{table}

\subsection{Kết nối Laptop, Desktop và Smartphone}
Máy bàn (Desktop) được nối bằng cáp Ethernet vào cổng LAN và cũng đã được thử
kết nối bằng Wi-Fi để chứng minh card không dây hoạt động. Laptop và điện thoại
được kết nối lần lượt vào mạng của nhóm; với thiết bị di động và Windows cần tắt
tính năng \emph{địa chỉ MAC ngẫu nhiên} cho mạng này, nếu không địa chỉ MAC hiện
trên router sẽ khác MAC thật và các luật lọc MAC ở trên sẽ không chạy đúng.

Trên mỗi thiết bị, sau khi kết nối cần kiểm tra: địa chỉ IP nhận được, cổng mặc
định (phải là \texttt{192.168.1.1}), băng tần và kênh đang dùng, kiểu bảo mật, và
khả năng liên lạc với router cũng như với thiết bị khác trong mạng. Kết quả đo
cụ thể được trình bày ở Mục~\ref{sec:experiment}.

\subsection{Thiết lập và kiểm chứng mạng Ad-hoc (IBSS) trên nền Linux}
\label{sec:adhoc-config}

Trước khi cấu hình, cần xác nhận nền tảng có hỗ trợ chế độ ngang hàng hay không.
Trên Windows~11, lệnh \texttt{netsh wlan show drivers} trả về dòng
\emph{Hosted network supported: No}, nghĩa là hệ điều hành không còn hỗ trợ chế
độ Ad-hoc (minh chứng tại Hình~\ref{fig:win-hosted}); Android và macOS cũng
không hỗ trợ. Vì vậy quy trình dưới đây được thực hiện trên Linux, sử dụng module
\texttt{mac80211\_hwsim} để mô phỏng phần cứng Wi-Fi và hai network namespace để
tách hai nút thành hai máy độc lập. Không được gọi nhầm chế độ hotspot hay
Wi-Fi Direct là \ac{ibss}.

\begin{enumerate}[leftmargin=*]
  \item \textbf{Kiểm tra nền tảng và tạo card Wi-Fi ảo.} Nạp module mô phỏng với
    hai radio: \texttt{sudo modprobe mac80211\_hwsim radios=2}, sau đó kiểm tra
    bằng \texttt{iw dev}. Mỗi radio tương ứng một giao diện (\texttt{wlan0},
    \texttt{wlan1}) có kênh, công suất phát và thuộc tính phần cứng riêng.
    Giao diện \texttt{P2P-device} do trình điều khiển tự tạo không ảnh hưởng
    đến kết nối Ad-hoc.

  \item \textbf{Tách mỗi radio vào một network namespace.} Nếu hai giao diện nằm
    chung một hệ thống, chúng dùng chung bảng định tuyến nên Linux coi cả hai
    địa chỉ IP là của chính nó và chuyển gói qua loopback, khiến lưu lượng không
    đi qua ``sóng'' vô tuyến ảo. Thực hiện:
    \texttt{sudo ip netns add ns1} và \texttt{sudo ip netns add ns2};
    \texttt{sudo iw phy phy0 set netns name ns1};
    \texttt{sudo iw phy phy1 set netns name ns2}. Kiểm tra lại bằng
    \texttt{sudo ip netns exec ns1 iw dev} và
    \texttt{sudo ip netns exec ns2 iw dev}.

  \item \textbf{Chuyển giao diện sang chế độ \ac{ibss} và tham gia mạng.} Với
    mỗi nút, hạ giao diện xuống trước khi đổi chế độ (nếu không, nhân hệ điều
    hành báo lỗi \emph{Device or resource busy}). Với Node~1 trong \texttt{ns1}:
    \texttt{sudo ip netns exec ns1 ip link set wlan0 down};
    \texttt{sudo ip netns exec ns1 iw dev wlan0 set type ibss};
    \texttt{sudo ip netns exec ns1 ip link set wlan0 up};
    \texttt{sudo ip netns exec ns1 iw dev wlan0 ibss join N09-AdHoc 2437};
    \texttt{sudo ip netns exec ns1 ip addr add 10.0.0.1/24 dev wlan0}. Tham số
    \texttt{N09-AdHoc} là \ac{ssid} của mạng và mọi nút phải trùng;
    \texttt{2437} là tần số tính bằng MHz, tương ứng kênh~6 của băng tần
    2,4~GHz và cũng phải trùng giữa các nút. Node~2 thực hiện tương tự với
    \texttt{ns2}, \texttt{wlan1} và địa chỉ \texttt{10.0.0.2/24}. Mọi lệnh đều
    chạy trong namespace tương ứng thông qua
    \texttt{sudo ip netns exec <namespace>}.

  \item \textbf{Kiểm tra liên kết ở tầng 802.11.} Dùng
    \texttt{iw dev wlan0 info} để xác nhận \texttt{type IBSS}, \ac{ssid},
    kênh và tần số; dùng \texttt{iw dev wlan0 station dump} để xem danh sách nút
    lân cận đã phát hiện. Nếu thấy địa chỉ MAC của nút còn lại ở trạng thái
    \emph{authorized} và \emph{associated}, hai nút đã nhận ra nhau.

  \item \textbf{Kiểm tra liên lạc ở tầng mạng.} Chạy
    \texttt{ping -c 4 10.0.0.2} từ Node~1 và chiều ngược lại từ Node~2. Gói tin
    được đóng khung 802.11, truyền qua ``sóng ảo'' của \texttt{hwsim} tới giao
    diện của nút kia. Kết quả đạt yêu cầu khi cả hai chiều đều nhận đủ 4/4 phản
    hồi và tỉ lệ mất gói bằng 0\,\%.

  \item \textbf{Ghi nhận và dọn dẹp.} Lưu ảnh chụp màn hình từng bước vào phần
    minh chứng (Mục~\ref{sec:evidence}) và ghi lại tham số thực tế. Sau khi kết
    thúc, gỡ namespace và module mô phỏng để trả máy về trạng thái ban đầu.
\end{enumerate}

\subsection{Bảng ghi cấu hình thực tế}
\begin{center}
\begin{tabularx}{\linewidth}{X X X}
  \toprule
  Hạng mục & Giá trị thực tế & Minh chứng/ghi chú \\
  \midrule
  Mẫu thiết bị/bộ định tuyến & TP-Link Archer C20 v5, OpenWrt (LuCI), tên máy
    \texttt{N09-Router} & Bảng~\ref{tab:cauhinh-thucte}; sơ đồ
    Hình~\ref{fig:sodo-n09} \\
  SSID và bảo mật & \texttt{25CLC04-N09} (2,4~GHz) và \texttt{25CLC04-N09-5G}
    (5~GHz), WPA2-PSK/CCMP; mạng khách \texttt{25CLC04-N09-Guest} &
    Bảng~\ref{tab:cauhinh-thucte} \\
  LAN/DHCP & \texttt{192.168.1.1/24}, DHCP thời hạn 12~h, IP tĩnh theo MAC cho
    máy bàn (\texttt{192.168.1.50}) & Bảng~\ref{tab:dhcp-leases} \\
  Laptop/Desktop/Smartphone & 4 thiết bị: máy bàn nối dây, một laptop và hai
    điện thoại Android; cả ba thiết bị không dây đều liên kết \ac{bss} 5~GHz &
    Hình~\ref{fig:sodo-n09}, Bảng~\ref{tab:ping-crosscheck} \\
  Ad-hoc hoặc phương án thay thế & IBSS trên Ubuntu 22.04 LTS (máy ảo VMware),
    \texttt{mac80211\_hwsim radios=2}; \ac{ssid} \texttt{N09-AdHoc}, kênh~6
    (2437~MHz), IP \texttt{10.0.0.1/24} và \texttt{10.0.0.2/24}, mạng mở &
    Hình~\ref{fig:sodo-adhoc}, Hình~\ref{fig:hwsim-radios}--\ref{fig:adhoc-ping};
    Windows~11 không hỗ trợ (Hình~\ref{fig:win-hosted}) \\
  \bottomrule
\end{tabularx}
\end{center}

\endgroup
